\section{Two small Pell proof reductions}
\label{sec:pell}
These are the recommended first changes. Both preserve theorem statements and use already imported infrastructure. Neither adds a helper, an import, a stronger hypothesis, or new executable data.

\subsection{Feed the square identity directly to Pell classification}
The target is \src{SignedPellResidue}{signed_residue_cases}, especially the block beginning at baseline line 62. The preceding argument already proves
\[
 (A\,|V|)^2=(A^2-1)|y|^2+1
\]
as an equality of natural numbers. The current proof then creates another equality with explicit multiplication, calls \code{nlinarith}, uses \code{omega} to turn addition into truncated subtraction, and finally passes that fact to \code{Pell.eq_pell}.

The additive equality is exactly the premise of \code{Nat.sub_eq_of_eq_add'}. It already contains the information that makes the natural subtraction exact. Normalize squares and multiplication association at the point of use:
\begin{lean}
obtain ⟨n, hx, hy⟩ := Pell.eq_pell hA
    (x := A * V.natAbs) (y := y.natAbs) (by
  change (A * V.natAbs) * (A * V.natAbs) -
    (A * A - 1) * y.natAbs * y.natAbs = 1
  simpa only [pow_two, mul_assoc] using
    Nat.sub_eq_of_eq_add' hn)
\end{lean}

The apostrophe is significant: it selects the orientation of the addition equation used here. The explicit \code{change} exposes the discriminant hidden behind the imported \code{Pell.d} definition; restricted simplification alone does not unfold that interface. Both details are present in the successfully checked patch.

The change removes the local \code{hp} bridge and three physical source lines: the complete module goes from 102 to 99 lines. It leaves the absolute-value conversion that produces \code{hn}, the index-parity argument, and all subsequent sign/congruence cases intact. Those later stages are adjacent to earlier proposals and receive no new credit here.

A possible alternative was to import an existing \code{eq_pell_of_sq} wrapper from a heavier module. That route is not recommended: this short adapter already exposes the exact mathematical step, while the heavier dependency would outweigh the small convenience. The supplied patch adds no dependency.

\subsection{Use the integer Pell norm theorem where the goal is integer-valued}
The target is \src{ThirteenFreshSeed}{auxiliary_exists}. Two local goals assert an integer norm identity for Pell sequence values. The baseline starts from \code{χ_sq}, an identity over natural numbers, casts it with \code{congrArg}, proves or supplies a guard for natural subtraction, and normalizes the resulting integer expression.

Pinned Mathlib already provides the integer invariant \code{Pell.pell_eqz}. Its discriminant conversion is \code{Pell.dz_val}. These are available through the module's existing imports. For the first norm identity, the entire proof becomes:
\begin{lean}
have hp : f ^ 2 - delta * (yn hA t : ℤ) ^ 2 = 1 := by
  simpa only [f, delta, xz, yz, dz_val, az,
    pow_two, mul_assoc] using pell_eqz hA t
\end{lean}

For the second, use the same integer invariant, substitute the existing coordinate equation \code{hx}, and retain the short polynomial normalization needed for the chosen sign of \code{y}:
\begin{lean}
have hnorm : (S : ℤ) ^ 2 * V ^ 2 -
    ((S : ℤ) ^ 2 - 1) * y ^ 2 = 1 := by
  have h := pell_eqz hS R
  simp only [xz, yz, dz_val, az, hx] at h
  dsimp [y]
  nlinarith only [h]
\end{lean}

This removes four lines at the first site and three at the second. The complete module goes from 217 to 210 physical lines. The local \code{hSS} subtraction guard disappears because the direct integer theorem already packages the correct discriminant. Other guards and casts used elsewhere in the construction remain necessary and are not removed.

The useful principle here is specific: choose the invariant in the scalar domain of the target equation. It is not a new general casting framework. Earlier integer-norm proposals address a different obligation, such as avoiding rational square arguments or minimizing a fundamental unit. This change substitutes an existing theorem for two constructive sequence-norm bridges. The earlier fresh-seed replacement itself retains those bridges, which provides a particularly close counterfactual comparison.

\subsection{Expected effect and limits}
Taken together, these two patches remove ten physical lines with no helper overhead. They reduce intermediate facts and two uses of natural-subtraction cast machinery. No compile-time, memory, or object-size improvement has been established. The validation section distinguishes complete modified-module checks from a project rebuild, and records the dependencies used. Their small scope makes these the best initial adoption candidates even though their absolute line savings are modest.
